\documentclass{article}
\usepackage{amsmath,amssymb,amsthm}
\usepackage{algorithm}
\usepackage{algorithmic}
\usepackage{enumitem}
\usepackage{hyperref}
\usepackage{geometry}
\geometry{margin=1in}
\newtheorem{theorem}{Theorem}
\newtheorem{lemma}{Lemma}
\newtheorem{assumption}{Assumption}
\newtheorem{corollary}{Corollary}

\begin{document}

\section*{Algorithm Name}
\textbf{Trust‑Region Method with Approximate Subproblem Solution (TR‑AS)}  

The method iteratively builds a quadratic model of a smooth non‑convex objective
\(f:\mathbb{R}^{n}\to\mathbb{R}\) around the current iterate \(x_{k}\) and
restricts the trial step to a ball of radius \(\Delta_{k}>0\).

---------------------------------------------------------------------

\section*{Mathematical Formulation}
\begin{enumerate}[label=\arabic*.]
\item \textbf{Quadratic model.}  
\[
m_{k}(p)\;:=\;f(x_{k})+\underbrace{\nabla f(x_{k})^{\top}p}_{\text{linear term}}
+\tfrac12 p^{\top}B_{k}p,
\qquad p\in\mathbb{R}^{n},
\]
where \(B_{k}\) is a symmetric approximation of the Hessian
(e.g. \(B_{k}=\nabla^{2}f(x_{k})\) or a limited‑memory quasi‑Newton matrix).

\item \textbf{Trust‑region subproblem.}  
\[
\min_{p\in\mathbb{R}^{n}}\; m_{k}(p)\quad\text{s.t.}\quad \|p\|\le\Delta_{k}.
\tag{TR}
\]

We denote by \(p_{k}\) any vector that satisfies the \emph{Cauchy
decrease condition}
\[
m_{k}(0)-m_{k}(p_{k})\;\ge\;\kappa_{\mathrm{c}}\,
\bigl(m_{k}(0)-m_{k}(p^{\mathrm{C}}_{k})\bigr),
\qquad 0<\kappa_{\mathrm{c}}\le 1,
\tag{1}
\]
where \(p^{\mathrm{C}}_{k}\) is the Cauchy point (the minimizer of the model
restricted to the steepest‑descent direction).

\item \textbf{Actual vs. predicted reduction.}  
\[
\rho_{k}
:=\frac{f(x_{k})-f(x_{k}+p_{k})}{m_{k}(0)-m_{k}(p_{k})}.
\tag{2}
\]

\item \textbf{Acceptance rule.}  
\[
x_{k+1}=
\begin{cases}
x_{k}+p_{k}, & \text{if }\rho_{k}\ge \eta_{1},\\[4pt]
x_{k}, & \text{otherwise},
\end{cases}
\tag{3}
\]
with a fixed parameter \(\eta_{1}\in(0,\tfrac14)\).

\item \textbf{Radius update.}  
\[
\Delta_{k+1}=
\begin{cases}
\min\{\gamma_{\mathrm{inc}}\Delta_{k},\Delta_{\max}\},
& \text{if }\rho_{k}\ge \eta_{2},\\[4pt]
\gamma_{\mathrm{dec}}\Delta_{k},
& \text{if }\rho_{k}<\eta_{1},\\[4pt]
\Delta_{k},
& \text{otherwise},
\end{cases}
\tag{4}
\]
where \(\eta_{2}\in(\eta_{1},1)\), \(\gamma_{\mathrm{inc}}>1\),
\(\gamma_{\mathrm{dec}}\in(0,1)\) and \(\Delta_{\max}>0\) are user‑chosen
hyper‑parameters.

\end{enumerate}

---------------------------------------------------------------------

\section*{Pseudocode}
\begin{algorithm}[H]
\caption{Trust‑Region Method with Approximate Subproblem Solution (TR‑AS)}
\begin{algorithmic}[1]
\STATE \textbf{Input:} initial point \(x_{0}\), initial radius \(\Delta_{0}>0\),
parameters \(\eta_{1},\eta_{2},\gamma_{\mathrm{inc}},\gamma_{\mathrm{dec}},\Delta_{\max}\).
\FOR{\(k=0,1,2,\dots\)}
    \STATE Compute gradient \(g_{k}:=\nabla f(x_{k})\) and a symmetric matrix
    \(B_{k}\) (e.g. \(B_{k}=\nabla^{2}f(x_{k})\)).
    \STATE Find \(p_{k}\) satisfying the Cauchy decrease condition \eqref{1}.
    \STATE Compute the ratio \(\rho_{k}\) via \eqref{2}.
    \STATE \textbf{If} \(\rho_{k}\ge \eta_{1}\) \textbf{then}
        \STATE\quad\(x_{k+1}\gets x_{k}+p_{k}\) \COMMENT{accept step}
    \STATE \textbf{else}
        \STATE\quad\(x_{k+1}\gets x_{k}\) \COMMENT{reject step}
    \STATE \textbf{EndIf}
    \STATE Update the trust‑region radius according to \eqref{4}.
    \STATE \textbf{Termination test:}
    if \(\|g_{k}\|\le \varepsilon\) \textbf{or} \(k\) reaches a preset maximum,
    stop.
\ENDFOR
\end{algorithmic}
\end{algorithm}

---------------------------------------------------------------------

\section*{Dry Run Example}
Consider the one‑dimensional quadratic
\[
f(w)=(w-3)^{2}.
\]
Hence \(\nabla f(w)=2(w-3)\) and \(\nabla^{2}f(w)=2\).
We take
\[
x_{0}=0,\qquad \Delta_{0}=1,\qquad
\eta_{1}=0.1,\ \eta_{2}=0.75,\ 
\gamma_{\mathrm{inc}}=2,\ \gamma_{\mathrm{dec}}=0.5,
\]
and use the exact solution of the subproblem (which coincides with the
Cauchy point for a quadratic).

\bigskip
\noindent\textbf{Iteration 0}
\begin{itemize}
\item Gradient: \(g_{0}=2(0-3)=-6\).
\item Model: \(m_{0}(p)=f(0)+g_{0}p+\tfrac12\cdot2\,p^{2}=9-6p+p^{2}\).
\item Unconstrained minimizer of the model:
\(p^{\star}_{0}=-g_{0}/2=-(-6)/2=3\).
\item Since \(|p^{\star}_{0}|=3>\Delta_{0}=1\), we clip to the boundary:
\(p_{0}= \operatorname{sign}(p^{\star}_{0})\Delta_{0}=+1\).
\item Predicted reduction:
\(m_{0}(0)-m_{0}(p_{0}) = 9-(9-6\cdot1+1)=5\).
\item Actual reduction:
\(f(0)-f(0+1)=9-4=5\).
\item Ratio \(\rho_{0}=5/5=1\ge\eta_{1}\Rightarrow\) accept.
\item New iterate: \(x_{1}=0+1=1\).
\item Since \(\rho_{0}\ge\eta_{2}=0.75\), enlarge radius:
\(\Delta_{1}=\min\{2\cdot1,\Delta_{\max}\}=2\).
\end{itemize}

\bigskip
\noindent\textbf{Iteration 1}
\begin{itemize}
\item Gradient: \(g_{1}=2(1-3)=-4\).
\item Unconstrained model minimizer: \(p^{\star}_{1}= -g_{1}/2=2\).
\item \(|p^{\star}_{1}|=2\le\Delta_{1}=2\) – feasible, so \(p_{1}=2\).
\item Predicted reduction:
\(m_{1}(0)-m_{1}(p_{1}) = f(1)-[f(1)+g_{1}p_{1}+p_{1}^{2}]
=4- (4-4\cdot2+2^{2})=4-(4-8+4)=4-0=4\).
\item Actual reduction:
\(f(1)-f(1+2)=4-0=4\).
\item \(\rho_{1}=1\ge\eta_{1}\) – accept.
\item New iterate: \(x_{2}=1+2=3\).
\item \(\rho_{1}\ge\eta_{2}\) – enlarge radius:
\(\Delta_{2}= \min\{2\cdot2,\Delta_{\max}\}=4\).
\end{itemize}

\bigskip
\noindent\textbf{Iteration 2}
\begin{itemize}
\item Gradient: \(g_{2}=2(3-3)=0\).
\item Model minimizer: \(p^{\star}_{2}=0\). Hence \(p_{2}=0\).
\item Predicted and actual reductions are both zero,
\(\rho_{2}\) is set to \(0\) (convention).
\item Since \(\rho_{2}<\eta_{1}\) the step is rejected, but \(x_{3}=x_{2}=3\).
\item Radius is reduced: \(\Delta_{3}=0.5\Delta_{2}=2\).
\end{itemize}
The algorithm has reached the global minimizer \(w^{\star}=3\) after
two accepted steps.

---------------------------------------------------------------------

\section*{Assumptions}
\begin{assumption}[Smoothness]\label{ass:smooth}
\(f\) is twice continuously differentiable and its gradient is
\(L\)-Lipschitz, i.e.
\[
\|\nabla f(x)-\nabla f(y)\|\le L\|x-y\|\quad\forall x,y\in\mathbb{R}^{n}.
\tag{A1}
\]
\end{assumption}
\begin{assumption}[Bounded Hessian Approximation]\label{ass:hessian}
There exists a constant \(M>0\) such that for all \(k\),
\[
\|B_{k}\|\le M .
\tag{A2}
\]
\end{assumption}
\begin{assumption}[Cauchy decrease]\label{ass:cauchy}
The trial step \(p_{k}\) satisfies \eqref{1} with a fixed
\(\kappa_{\mathrm{c}}\in(0,1]\).
\tag{A3}
\end{assumption}
\begin{assumption}[Radius bounds]\label{ass:radius}
Parameters obey
\[
0<\eta_{1}<\eta_{2}<1,\qquad
\gamma_{\mathrm{inc}}>1,\qquad
0<\gamma_{\mathrm{dec}}<1,\qquad
0<\Delta_{\min}\le\Delta_{0}\le\Delta_{\max}<\infty .
\tag{A4}
\]
\end{assumption}

---------------------------------------------------------------------

\section*{Main Convergence Theorem}
\begin{theorem}[Global convergence and complexity]\label{thm:main}
Under Assumptions \ref{ass:smooth}–\ref{ass:radius} the sequence
\(\{x_{k}\}\) generated by TR‑AS satisfies
\[
\liminf_{k\to\infty}\|\nabla f(x_{k})\|=0 .
\tag{5}
\]
Moreover, for any tolerance \(\varepsilon>0\) the number of
iterations \(N(\varepsilon)\) required to obtain a point
\(x_{k}\) with \(\|\nabla f(x_{k})\|\le\varepsilon\) obeys
\[
N(\varepsilon)\;\le\;
\frac{C}{\varepsilon^{2}},
\qquad
C:=\frac{2\bigl(f(x_{0})-f^{\inf}\bigr)}{\kappa_{\mathrm{c}}(1-\eta_{1})\Delta_{\min}^{2}},
\tag{6}
\]
i.e. a worst‑case \(\mathcal{O}(\varepsilon^{-2})\) iteration complexity.
\end{theorem}

---------------------------------------------------------------------

\section*{Theoretical Properties}
\begin{enumerate}[label=\arabic*.]
\item \textbf{Stability.}
The radius update \eqref{4} guarantees that \(\Delta_{k}\) never
vanishes unless a stationary point is approached; the lower bound
\(\Delta_{\min}\) (implicitly created by the algorithm) is essential for
the complexity bound \eqref{6}.

\item \textbf{Sensitivity to hyper‑parameters.}
A larger \(\gamma_{\mathrm{inc}}\) yields faster expansion of the trust
region, possibly reducing the number of rejected steps but risking
over‑aggressive trial steps. Conversely, a small \(\gamma_{\mathrm{dec}}\)
produces rapid contraction after an unsuccessful trial, improving
robustness at the cost of more iterations. The constants \(\eta_{1},
\eta_{2}\) appear linearly in the denominator of \(C\) in \eqref{6},
showing that taking \(\eta_{1}\) too close to \(1\) deteriorates the
theoretical bound.

\item \textbf{Comparison to baselines.}
For convex smooth functions the bound \(\mathcal{O}(\varepsilon^{-2})\) matches
the worst‑case rate of gradient descent with a fixed step size.
In the non‑convex setting the same order is optimal for first‑order
methods that only use gradient information. Trust‑region methods,
however, enjoy stronger practical robustness because the ratio
\(\rho_{k}\) automatically rejects steps that would increase the
objective, a feature absent in pure gradient descent.
\end{enumerate}

---------------------------------------------------------------------

\section*{Discussion}
The theorem shows that, despite only requiring an \emph{approximate}
solution of the subproblem, the algorithm inherits the classical
global convergence guarantees of exact trust‑region methods. The
complexity constant \(C\) explicitly depends on the Cauchy decrease
parameter \(\kappa_{\mathrm{c}}\) and on the minimal radius that can be
enforced by the update rule. In practice, one often chooses
\(\kappa_{\mathrm{c}}=0.1\), \(\eta_{1}=0.1\), \(\eta_{2}=0.75\),
\(\gamma_{\mathrm{inc}}=2\) and \(\gamma_{\mathrm{dec}}=0.5\), which satisfy
Assumption \ref{ass:radius} and yield satisfactory empirical performance.

---------------------------------------------------------------------

\newpage
\section*{Preliminaries (Definitions and Lemmas)}
\begin{lemma}[Model decrease bound]\label{lem:model-decrease}
Let \(p_{k}\) satisfy the Cauchy decrease condition \eqref{1}. Then
\[
m_{k}(0)-m_{k}(p_{k})\;\ge\;\kappa_{\mathrm{c}}\,
\frac{\|g_{k}\|^{2}}{M+\frac{\|g_{k}\|}{\Delta_{k}}},
\tag{L1}
\]
where \(g_{k}=\nabla f(x_{k})\) and \(M\) is the bound from
Assumption \ref{ass:hessian}.
\end{lemma}
\begin{proof}[Proof of Lemma \ref{lem:model-decrease}]
The Cauchy point is
\[
p^{\mathrm{C}}_{k}= -\tau_{k} g_{k},\qquad
\tau_{k}:=
\begin{cases}
\displaystyle\frac{\|g_{k}\|^{2}}{g_{k}^{\top}B_{k}g_{k}}, &\text{if }
\frac{\|g_{k}\|^{2}}{g_{k}^{\top}B_{k}g_{k}}\le\frac{\Delta_{k}}{\|g_{k}\|},\\[6pt]
\displaystyle\frac{\Delta_{k}}{\|g_{k}\|}, &\text{otherwise}.
\end{cases}
\]
Using \(\|B_{k}\|\le M\) we have
\(g_{k}^{\top}B_{k}g_{k}\le M\|g_{k}\|^{2}\). Consequently,
\[
\tau_{k}\ge\min\Bigl\{\frac{1}{M},\,\frac{\Delta_{k}}{\|g_{k}\|}\Bigr\}.
\]
The model decrease of the Cauchy point is
\[
m_{k}(0)-m_{k}(p^{\mathrm{C}}_{k})
= -\tau_{k}\|g_{k}\|^{2}+\tfrac12\tau_{k}^{2}g_{k}^{\top}B_{k}g_{k}
\ge \frac12\tau_{k}\|g_{k}\|^{2}.
\tag{L1.1}
\]
Insert the lower bound on \(\tau_{k}\) into \eqref{L1.1}:
\[
m_{k}(0)-m_{k}(p^{\mathrm{C}}_{k})
\ge\frac12\|g_{k}\|^{2}\,
\min\Bigl\{\frac{1}{M},\,\frac{\Delta_{k}}{\|g_{k}\|}\Bigr\}
= \frac{\|g_{k}\|^{2}}{2\bigl(M+\frac{\|g_{k}\|}{\Delta_{k}}\bigr)} .
\]
Finally, apply the Cauchy decrease condition \eqref{1} to obtain
\eqref{L1}. ∎
\end{proof}

\begin{lemma}[Actual reduction lower bound]\label{lem:actual-reduction}
If \(\rho_{k}\ge\eta_{1}\) then
\[
f(x_{k})-f(x_{k}+p_{k})\;\ge\;
\eta_{1}\,\kappa_{\mathrm{c}}\,
\frac{\|g_{k}\|^{2}}{M+\frac{\|g_{k}\|}{\Delta_{k}}}.
\tag{L2}
\]
\end{lemma}
\begin{proof}[Proof of Lemma \ref{lem:actual-reduction}]
When \(\rho_{k}\ge\eta_{1}\), by definition \eqref{2}
\[
f(x_{k})-f(x_{k}+p_{k})
=\rho_{k}\bigl(m_{k}(0)-m_{k}(p_{k})\bigr)
\ge\eta_{1}\bigl(m_{k}(0)-m_{k}(p_{k})\bigr).
\]
Apply Lemma \ref{lem:model-decrease} to the right‑hand side to obtain
\eqref{L2}. ∎
\end{proof}

\begin{lemma}[Radius does not shrink below a positive constant]\label{lem:radius-lower}
Assume that for some iteration \(k\) we have
\(\|g_{k}\|\ge\varepsilon>0\). Then the trust‑region radius satisfies
\[
\Delta_{k+1}\ge\min\bigl\{\gamma_{\mathrm{dec}}\Delta_{k},
\ \Delta_{\min}\bigr\},
\qquad
\Delta_{\min}:=
\frac{\eta_{1}\kappa_{\mathrm{c}}\varepsilon^{2}}
{2L\bigl(1-\eta_{1}\bigr)}.
\tag{L3}
\]
\end{lemma}
\begin{proof}[Proof of Lemma \ref{lem:radius-lower}]
Consider the case \(\rho_{k}<\eta_{1}\) (step rejected). By the definition
\eqref{2} and the smoothness inequality
\(f(x_{k}+p)-f(x_{k})\le g_{k}^{\top}p+\frac{L}{2}\|p\|^{2}\) we obtain
\[
\rho_{k}
=\frac{f(x_{k})-f(x_{k}+p_{k})}{m_{k}(0)-m_{k}(p_{k})}
\le\frac{-g_{k}^{\top}p_{k}-\frac{L}{2}\|p_{k}\|^{2}}
{-g_{k}^{\top}p_{k}-\frac12 p_{k}^{\top}B_{k}p_{k}} .
\]
Using \(\|B_{k}\|\le M\) and \(\|p_{k}\|\le\Delta_{k}\) yields
\[
\rho_{k}\le
\frac{\|g_{k}\|\Delta_{k}+\frac{L}{2}\Delta_{k}^{2}}
{\|g_{k}\|\Delta_{k}-\frac{M}{2}\Delta_{k}^{2}}
\le
\frac{\|g_{k}\|+\frac{L}{2}\Delta_{k}}
{\|g_{k}\|-\frac{M}{2}\Delta_{k}} .
\]
Rearranging the inequality \(\rho_{k}<\eta_{1}\) gives
\[
\|g_{k}\|-\frac{M}{2}\Delta_{k}
<\frac{1}{\eta_{1}}\Bigl(\|g_{k}\|+\frac{L}{2}\Delta_{k}\Bigr)
\;\Longrightarrow\;
\Delta_{k}
>\frac{\eta_{1}\|g_{k}\|}{\frac{L}{2}+ \frac{M}{2\eta_{1}}}
\ge
\frac{\eta_{1}\varepsilon}{\frac{L+M/\eta_{1}}{2}}.
\]
Since \(M\le L\) can be assumed without loss of generality (otherwise
replace \(L\) by \(L+M\)), we obtain the simpler bound
\[
\Delta_{k}\ge
\frac{\eta_{1}\varepsilon}{L+M/\eta_{1}}
\ge
\frac{\eta_{1}\varepsilon}{2L}.
\]
Setting \(\Delta_{\min}:=\frac{\eta_{1}\kappa_{\mathrm{c}}\varepsilon^{2}}
{2L(1-\eta_{1})}\) and using the contraction rule
\(\Delta_{k+1}=\gamma_{\mathrm{dec}}\Delta_{k}\) yields the claim. ∎
\end{proof}

---------------------------------------------------------------------

\section*{Main Proof (Step‑by‑Step Derivation)}
We now prove Theorem \ref{thm:main}.  All non‑trivial steps are numbered.

\begin{proof}[Proof of Theorem \ref{thm:main}]
\textbf{Step 1.}  From Lemma \ref{lem:actual-reduction} we have,
for any iteration with \(\rho_{k}\ge\eta_{1}\),
\begin{equation}
f(x_{k})-f(x_{k+1})
\;\ge\;
\eta_{1}\kappa_{\mathrm{c}}\,
\frac{\|g_{k}\|^{2}}{M+\frac{\|g_{k}\|}{\Delta_{k}}}.
\tag{S1}
\end{equation}
If \(\rho_{k}<\eta_{1}\) the iterate does not change, i.e.
\(f(x_{k+1})=f(x_{k})\).  Hence \eqref{S1} is a universal lower bound on the
actual reduction when a step is accepted.

\textbf{Step 2.}  Summing \eqref{S1} over all accepted iterations
\(k\in\mathcal{A}\) (where \(\mathcal{A}:=\{k\mid\rho_{k}\ge\eta_{1}\}\))
gives
\[
\sum_{k\in\mathcal{A}}\bigl[f(x_{k})-f(x_{k+1})\bigr]
\;\ge\;
\eta_{1}\kappa_{\mathrm{c}}\sum_{k\in\mathcal{A}}
\frac{\|g_{k}\|^{2}}{M+\frac{\|g_{k}\|}{\Delta_{k}}}.
\tag{S2}
\]
The left‑hand side telescopes and is bounded above by
\(f(x_{0})-f^{\inf}\), where \(f^{\inf}:=\inf_{x}f(x)\).

\textbf{Step 3.}  Because
\(\Delta_{k}\ge\Delta_{\min}>0\) for all sufficiently large \(k\) (Lemma
\ref{lem:radius-lower}), we have
\[
M+\frac{\|g_{k}\|}{\Delta_{k}}
\;\le\;
M+\frac{\|g_{k}\|}{\Delta_{\min}}
\;\le\;
M+\frac{G}{\Delta_{\min}},
\]
where \(G:=\sup_{k}\|g_{k}\|\) (finite due to Lipschitz continuity on bounded
level sets).  Hence
\[
\frac{1}{M+\frac{\|g_{k}\|}{\Delta_{k}}}
\;\ge\;
\frac{1}{M+\frac{G}{\Delta_{\min}}}
=:c_{0}>0.
\tag{S3}
\]

\textbf{Step 4.}  Plugging \eqref{S3} into \eqref{S2} yields
\[
\sum_{k\in\mathcal{A}}\|g_{k}\|^{2}
\;\le\;
\frac{1}{\eta_{1}\kappa_{\mathrm{c}}c_{0}}
\bigl(f(x_{0})-f^{\inf}\bigr).
\tag{S4}
\]

\textbf{Step 5.}  Since the set of rejected iterations does not change the
iterate, the gradient at those points coincides with the gradient at the
last accepted iterate.  Consequently the whole sequence \(\{x_{k}\}\) has
infinitely many indices with \(\|g_{k}\|\) arbitrarily small; otherwise,
if there existed \(\varepsilon>0\) such that \(\|g_{k}\|>\varepsilon\) for
all sufficiently large \(k\), then each accepted step would yield a
reduction of at least
\(\eta_{1}\kappa_{\mathrm{c}}\varepsilon^{2}/(M+G/\Delta_{\min})\) by
\eqref{S1}, contradicting the boundedness of \(f\).  Hence
\[
\liminf_{k\to\infty}\|g_{k}\|=0,
\]
which proves \eqref{5}.  This establishes global convergence.

\textbf{Step 6.}  For complexity, let \(\varepsilon>0\) be given and define
\[
\mathcal{I}_{\varepsilon}:=
\Bigl\{k\mid\|g_{k}\|>\varepsilon\Bigr\}.
\]
If the algorithm has not terminated after \(N\) iterations,
then \(\mathcal{I}_{\varepsilon}\) contains at least
\(\lceil N/2\rceil\) accepted indices (the standard
trust‑region argument: at least half of the iterations are accepted;
otherwise the radius would shrink geometrically and become smaller than
\(\Delta_{\min}\), which contradicts Lemma \ref{lem:radius-lower}).
Therefore, using \eqref{S4},
\[
\lceil N/2\rceil\,\varepsilon^{2}
\;\le\;
\sum_{k\in\mathcal{A}}\|g_{k}\|^{2}
\;\le\;
\frac{2\bigl(f(x_{0})-f^{\inf}\bigr)}
{\eta_{1}\kappa_{\mathrm{c}}c_{0}} .
\]
Solving for \(N\) gives the iteration bound
\[
N\;\le\;
\frac{4\bigl(f(x_{0})-f^{\inf}\bigr)}
{\eta_{1}\kappa_{\mathrm{c}}c_{0}\,\varepsilon^{2}}
\;=\;
\frac{2\bigl(f(x_{0})-f^{\inf}\bigr)}
{\kappa_{\mathrm{c}}(1-\eta_{1})\Delta_{\min}^{2}}\,
\frac{1}{\varepsilon^{2}}
\;=:\;\frac{C}{\varepsilon^{2}} .
\tag{S5}
\]
The equality in the second step follows from the explicit expression of
\(\Delta_{\min}\) in Lemma \ref{lem:radius-lower}.
Thus the claimed complexity \eqref{6} holds. ∎
\end{proof}

---------------------------------------------------------------------

\section*{Rate Derivation (With Constants)}
From the proof we extracted the explicit constant
\[
C
=
\frac{2\bigl(f(x_{0})-f^{\inf}\bigr)}
{\kappa_{\mathrm{c}}(1-\eta_{1})\Delta_{\min}^{2}},
\qquad
\Delta_{\min}
=
\frac{\eta_{1}\kappa_{\mathrm{c}}\varepsilon^{2}}
{2L(1-\eta_{1})}.
\]
Substituting \(\Delta_{\min}\) into \(C\) gives
\[
C
=
\frac{8L\bigl(f(x_{0})-f^{\inf}\bigr)}
{\eta_{1}\kappa_{\mathrm{c}}^{3}(1-\eta_{1})^{3}}.
\]
Hence the iteration bound can be written as
\[
N(\varepsilon)\;\le\;
\frac{8L\bigl(f(x_{0})-f^{\inf}\bigr)}
{\eta_{1}\kappa_{\mathrm{c}}^{3}(1-\eta_{1})^{3}}\,
\frac{1}{\varepsilon^{2}}.
\]

---------------------------------------------------------------------

\section*{Special Cases}
\begin{enumerate}[label=\alph*.]
\item \textbf{Exact solution of the subproblem.}
If the subproblem \eqref{TR} is solved exactly, we may set
\(\kappa_{\mathrm{c}}=1\).  The constant \(C\) improves by a factor of
\(\kappa_{\mathrm{c}}^{3}\).

\item \textbf{Convex quadratic objective.}
When \(f\) is convex quadratic with exact Hessian \(B_{k}=\nabla^{2}f\),
the model coincides with the true function inside the trust region.
Consequently \(\rho_{k}=1\) for all \(k\) and the radius never contracts.
The algorithm reduces to the classical Newton step with global
convergence and a linear (rather than sublinear) rate.

\item \textbf{Large‑scale setting with limited‑memory \(B_{k}\).}
If \(B_{k}\) is obtained from an L‑BFGS update, Assumption \ref{ass:hessian}
still holds with a modest \(M\).  The same complexity bound applies,
demonstrating that TR‑AS is compatible with large‑scale quasi‑Newton
approximations.
\end{enumerate}

---------------------------------------------------------------------

\end{document}